\begin{ouverture}[ensembles]{2}{Réels et intervalles}{Partie 1 · Nombres et calculs}
{Ensembles de nombres · Encadrements · Distances}
\objectif{Distinguer valeur exacte et valeur approchée.}
\objectif{Lire et représenter les intervalles.}
\objectif{Traduire une distance sur la droite numérique.}
\prerequis{Fractions, nombres relatifs et ordre.}
\end{ouverture}
\pagestyle{otMain}\essentiel\fontsize{11}{15}\selectfont
\essSection{1}{Les nombres et leurs écritures}
\begin{definition}
$\mathbb N$ contient les entiers naturels, dont 0 ; $\mathbb Z$ les entiers relatifs ; $\mathbb D$ les nombres ayant une écriture décimale finie ; $\mathbb Q$ les quotients d’entiers de dénominateur non nul. Tous appartiennent à $\mathbb R$.
\[\mathbb N\subset\mathbb Z\subset\mathbb D\subset\mathbb Q\subset\mathbb R.\]
Un irrationnel est un réel qui n’appartient pas à $\mathbb Q$, par exemple $\sqrt2$ ou $\pi$.
\end{definition}
\begin{exemple}$3/8=0{,}375$ est décimal ; $1/3$ est rationnel mais non décimal. $\sqrt{25}=5$ est entier : une racine carrée n’est pas toujours irrationnelle.\end{exemple}
\begin{savoirfaire}[Encadrer et arrondir]
Un encadrement $a\le x\le b$ a pour amplitude $b-a$. Pour arrondir au centième, repérer les deux centièmes consécutifs qui encadrent le nombre puis choisir le plus proche. Préciser l’unité et la précision dans un problème de mesure.
\end{savoirfaire}
\begin{attention}Un affichage de calculatrice est souvent approché. $1{,}414$ n’est pas égal à $\sqrt2$ ; écrire $\sqrt2\approx1{,}414$.\end{attention}
\clearpage
\essSection{2}{Intervalles et distance}
\begin{definition}
$[a;b]$ contient les réels tels que $a\le x\le b$ ; $]a;b[$ impose $a<x<b$. Un crochet tourné vers le nombre signifie que la borne est incluse. Aux infinis, le crochet est toujours ouvert. Le programme officiel écrit aussi $[a,b]$ ou $[a,+\infty[$ avec une virgule ; ce livre utilise le point-virgule pour éviter toute confusion avec la virgule décimale, comme dans $[1{,}5;2]$.
\end{definition}
\begin{exemple}$x>-2$ s’écrit $x\in]-2;+\infty[$. La réunion $I\cup J$ signifie « dans $I$ ou dans $J$ » ; l’intersection $I\cap J$, « dans les deux ».\end{exemple}
\begin{definition}La valeur absolue $|a|$ est la distance entre $a$ et 0. La distance entre $a$ et $b$ est $|a-b|$. Pour $r\ge0$ :
\[|x-a|\le r\iff a-r\le x\le a+r.\]
\end{definition}
\begin{savoirfaire}[Lire un intervalle centré]
$|x-3|\le2$ signifie que $x$ est à distance au plus 2 de 3, donc $x\in[1;5]$. Le centre est 3 et le rayon est 2.
\end{savoirfaire}
\begin{attention}Une longueur est positive ; une abscisse peut être négative. Pour une borne exclue, ne pas remplacer une inégalité stricte par une inégalité large.\end{attention}
\clearpage\rubrique{Découverte}
\begin{blocLivre}{ex-2-1}
\begin{exercice}[1]{Le plus petit ensemble}
Pour chacun des nombres $-4$, $7/8$, $2/3$, $\sqrt{49}$ et $\sqrt3$, donner le plus petit ensemble parmi $\mathbb N,\mathbb Z,\mathbb D,\mathbb Q,\mathbb R$ qui le contient.
\end{exercice}
\end{blocLivre}
\begin{blocLivre}{ex-2-2}
\begin{exercice}[1]{Comparer exactement}
Ranger $-3/2$, $-1{,}4$, $7/5$ et $\sqrt2$ dans l’ordre croissant sans remplacer la racine par un arrondi grossier.
\end{exercice}
\end{blocLivre}
\begin{blocLivre}{ex-2-3}
\begin{exercice}[1]{Du texte à l’intervalle}
\begin{questions}
\item Écrire sous forme d’intervalle : $-2\le x<5$ ; $x\ge3$ ; $x<-1$.
\item Dire si 5 appartient au premier intervalle.
\end{questions}
\end{exercice}
\end{blocLivre}
\begin{blocLivre}{ex-2-4}
\begin{exercice}[1]{Union et intersection}
Pour $I=[-3;2]$ et $J=]0;5]$, déterminer $I\cap J$ et $I\cup J$. Représenter les deux intervalles sur une droite.
\label{fig-ex-2-4}
\figureGeo{intervalles}
\end{exercice}
\end{blocLivre}
\begin{blocLivre}{ex-2-5}
\begin{exercice}[1]{Distances sur une droite}
Calculer la distance entre les abscisses $-4$ et 3, puis entre $2{,}5$ et $-1{,}5$. Écrire chaque calcul avec une valeur absolue.
\end{exercice}
\end{blocLivre}

\clearpage\rubrique{Entraînement}
\begin{blocLivre}{ex-2-6}
\begin{exercice}[2]{Un intervalle centré}
Résoudre $|x-2|\le3$ et représenter l’ensemble obtenu.
\label{fig-ex-2-6}
\figureGeo{intervalle-centre}
\end{exercice}
\end{blocLivre}
\begin{blocLivre}{ex-2-7}
\begin{exercice}[2]{Un encadrement de racine}
Démontrer $2{,}23<\sqrt5<2{,}24$. En déduire l’arrondi de $\sqrt5$ au dixième.
\end{exercice}
\end{blocLivre}
\begin{blocLivre}{ex-2-8}
\begin{exercice}[2]{Une erreur absolue}
Un capteur annonce $12{,}4$ cm avec une erreur absolue d’au plus $0{,}15$ cm. Donner l’intervalle des longueurs possibles et sa largeur.
\end{exercice}
\end{blocLivre}
\begin{blocLivre}{ex-2-9}
\begin{exercice}[2]{Complémentaire d’un intervalle}
Dans $\mathbb R$, décrire les nombres qui n’appartiennent pas à $]-2;4]$. Donner une écriture par réunion.
\end{exercice}
\end{blocLivre}
\begin{blocLivre}{ex-2-10}
\begin{exercice}[2]{Mesures et chiffres significatifs}
Un rectangle mesure $12{,}3$ cm sur $4{,}56$ cm. Calculer son aire, puis l’arrondir à trois chiffres significatifs. Convertir la valeur non arrondie en mètres carrés. Effectuer aussi ces conversions : 750 mL en L ; $2{,}5$ h en heures et minutes ; 72 km/h en m/s ; $1{,}4$ kg en g ; 3 dm$^3$ en L.
\end{exercice}
\end{blocLivre}

\clearpage\rubrique{Problèmes}
\begin{blocLivre}{ex-2-11}
\begin{exercice}[3]{Deux contraintes de mesure}
Une longueur $x$ en centimètres vérifie $|x-10|\le2$ et $9<x\le13$. Quel est l’ensemble exact des valeurs possibles ?
\end{exercice}
\end{blocLivre}
\begin{blocLivre}{ex-2-12}
\begin{exercice}[3]{Un décimal caché}
Parmi $13/40$, $7/12$ et $21/28$, quels nombres sont décimaux ? On pourra chercher une écriture de dénominateur puissance de 10.
\end{exercice}
\end{blocLivre}
\begin{blocLivre}{ex-2-13}
\begin{exercice}[3]{Un faux calcul sur les racines}
Une élève affirme $\sqrt{a+b}=\sqrt a+\sqrt b$ pour tous $a,b\ge0$. Réfuter l’affirmation, puis comparer les deux membres pour $a=b=4$.
\end{exercice}
\end{blocLivre}
\begin{blocLivre}{ex-2-14}
\begin{exercice}[3]{Une amplitude imposée}
On sait $1{,}732<\sqrt3<1{,}733$. Encadrer $2+3\sqrt3$ avec une amplitude de $0{,}003$. Donner aussi son arrondi au centième.
\end{exercice}
\end{blocLivre}
\begin{blocLivre}{ex-2-15}
\begin{exercice}[3]{Balayer pour encadrer}
Pour encadrer $\sqrt2$ au centième, on teste les entiers $k$ et on compare $k^2$ à 20000. Expliquer pourquoi, puis trouver les deux centièmes consécutifs convenables.
\end{exercice}
\end{blocLivre}

\clearpage
